% !TEX root = ../../main.tex
% C2.1 — Các hệ thống liên quan hiện có trên thị trường (RÚT GỌN 2026-09-29; bản cũ ở report/archive/)
% Mỗi hệ thống: nhà phát triển, các hình thức tìm kiếm, triển khai, liên hệ với đề tài. BẮT BUỘC có nguồn. Chưa dùng thuật ngữ ReID/TBPS (định nghĩa ở Chương 3).
\section{Các hệ thống liên quan hiện có trên thị trường}
\label{sec:2.1}

Mục này khảo sát ba sản phẩm thương mại đại diện cho ba hướng tiếp cận tìm kiếm người trong dữ liệu camera: phân tích video toàn diện (BriefCam), tìm theo ngoại hình (Avigilon) và tìm bằng câu mô tả trên nền tảng đám mây (Verkada); các công trình học thuật được trình bày tại Chương~\ref{chapter:lythuyet}.

\subsection{BriefCam}
\label{sec:2.1.1}

BriefCam là phần mềm phân tích nội dung video ứng dụng học sâu, được Canon mua lại năm 2018~\cite{briefcam2018canon} và hiện được phân phối như một sản phẩm của Milestone Systems~\cite{milestone2026briefcam}. Chức năng được chia thành ba module: \emph{Review} phục vụ điều tra trên video đã ghi, cho phép lọc người và phương tiện theo thuộc tính như màu sắc, kích thước, tốc độ, hướng di chuyển, tìm đối tượng có ngoại hình tương tự một đối tượng mẫu, và tóm tắt hàng giờ video thành một đoạn ngắn (Video Synopsis); \emph{Respond} phát cảnh báo thời gian thực theo quy tắc; và \emph{Research} cung cấp thống kê, bản đồ nhiệt và đếm người~\cite{milestone2026briefcam}. Hình~\ref{fig:briefcam} minh họa giao diện tìm kiếm của sản phẩm. BriefCam có thể cài đặt tại chỗ hoặc trên đám mây, nhưng máy chủ xử lý bắt buộc phải có card đồ họa NVIDIA~\cite{milestone2026briefcamfaq}.

\begin{figure}[H]
    \centering
    \includegraphics[width=0.6\textwidth]{briefcam.png}
    \caption{Giao diện tìm kiếm và lọc đối tượng của BriefCam~\cite{milestone2026briefcam}}
    \label{fig:briefcam}
\end{figure}

\paragraph{Liên hệ với đề tài.} BriefCam hỗ trợ tìm theo thuộc tính và theo mẫu ngoại hình, trong đó mẫu là một đối tượng chọn trong video đã ghi thay vì một ảnh tải lên từ bên ngoài như trong đề tài; tài liệu công bố không đề cập tìm kiếm bằng câu mô tả tự nhiên.

\subsection{Avigilon Appearance Search}
\label{sec:2.1.2}

Avigilon Appearance Search là công nghệ tìm kiếm video dựa trên học sâu của Avigilon, thuộc Motorola Solutions, dùng để tìm một người hay phương tiện trong nhiều giờ video trên nhiều camera~\cite{motorola2021appearance}. Người dùng có thể bắt đầu tìm kiếm bằng cách chọn đặc điểm ngoại hình từ danh sách có sẵn như màu quần áo, giới tính, nhóm tuổi; bằng cách tải lên một bức ảnh; hoặc bằng cách chọn một đối tượng trong video làm mẫu~\cite{motorola2021appearance}. Kết quả giúp lần theo lộ trình của đối tượng và có thể được đánh dấu, trích xuất thành bộ bằng chứng về diễn biến sự việc. Hình~\ref{fig:avigilon} minh họa giao diện khi đánh dấu một người. Công nghệ được tích hợp trong phần mềm quản lý video Avigilon Control Center, chạy trên camera và đầu ghi của hãng hoặc trên thiết bị Avigilon AI Appliance cho camera ONVIF của hãng khác.

\begin{figure}[H]
    \centering
    \includegraphics[width=0.45\textwidth]{avigilon.png}
    \caption{Giao diện phần mềm Avigilon Control Center với đối tượng được đánh dấu~\cite{avigilon2026appearance}}
    \label{fig:avigilon}
\end{figure}

\paragraph{Liên hệ với đề tài.} Avigilon hỗ trợ hai hình thức mô tả của đề tài là chọn thuộc tính và dùng ảnh tham chiếu, và cho phép tổng hợp kết quả thành bằng chứng, tương tự hồ sơ vụ việc của đề tài; tài liệu không đề cập tìm bằng câu mô tả tự do, và hệ thống gắn với phần mềm, thiết bị của hãng.

\subsection{Verkada AI-Powered Search}
\label{sec:2.1.3}

Verkada AI-Powered Search, giới thiệu năm 2024 trên nền tảng quản lý an ninh đám mây Command, cho phép tìm người và phương tiện bằng ngôn ngữ đời thường, ví dụ ``person wearing helmet'', thay vì chọn từ danh sách thuộc tính~\cite{verkada2024aisearch}, và tìm bằng ảnh tải lên của người, phương tiện hay đồ vật~\cite{verkada2024reverse}. Theo bài viết kỹ thuật của hãng~\cite{verkada2025scaling}, camera tạo các ảnh cắt của đối tượng; trên đám mây, ảnh cắt được mã hóa thành vector bằng mô hình CLIP, vốn đưa ảnh và văn bản vào cùng một không gian biểu diễn, và lưu vào một cơ sở dữ liệu vector; câu truy vấn được mã hóa theo cùng cách để tìm các ảnh gần nhất. Hình~\ref{fig:verkada} minh họa kết quả của một truy vấn bằng câu mô tả. Hệ thống vận hành tốt nhất với camera của Verkada; camera hãng khác kết nối qua thiết bị Command Connector với độ trễ cao hơn~\cite{verkada2024connector}.

\begin{figure}[H]
    \centering
    \includegraphics[width=0.52\textwidth]{verkada.png}
    \caption{Kết quả tìm kiếm bằng câu mô tả trên nền tảng Verkada Command~\cite{verkada2024aisearch}}
    \label{fig:verkada}
\end{figure}

\paragraph{Liên hệ với đề tài.} Verkada gần với đề tài nhất về kỹ thuật: cùng hỗ trợ tìm bằng câu mô tả và bằng ảnh dựa trên không gian biểu diễn chung và cơ sở dữ liệu vector. Khác biệt nằm ở mô hình triển khai: Verkada xử lý trên hạ tầng đám mây của hãng, còn đề tài triển khai tại chỗ trên máy tính phổ thông.
